\documentclass[10pt,a4paper]{amsart}
\usepackage[margin=19mm]{geometry}
\usepackage{amsmath,amssymb,mathtools}
\usepackage[scr=rsfs,cal=euler]{mathalfa}
\usepackage{xfrac}
\usepackage[unicode,hidelinks,bookmarks=false]{hyperref}
\newcommand{\ie}{i.e.\ }
\newcommand{\cf}{cf.\ }
\newcommand{\resp}{respectively\ }
\newcommand{\Subsection}{\subsection}
\providecommand{\nobreakdash}{\nobreak\hbox{-}}
\setlength{\emergencystretch}{2em}
\multlinegap0pt
\usepackage{bm}
\usepackage{bbm}
\usepackage{dsfont}
\usepackage{xspace}
\usepackage{cancel}
\usepackage{mathtools}
\usepackage{amsthm}
\makeatletter
\@ifundefined{theorem}{\newtheorem{theorem}{Theorem}}{}
\makeatother
\title[On Posets of Classes of Automorphic Subgroups of Finite Grou]{On Posets of Classes of Automorphic Subgroups of Finite Groups}
\author{Sachin Ballal, Tushar Halder}
\date{Source: arXiv:2511.00433v1 (CC BY 4.0)}
\begin{document}
\maketitle

\noindent\emph{Source and licence.} On Posets of Classes of Automorphic Subgroups of Finite Groups. Version arXiv:2511.00433v1 (CC BY 4.0), distributed under the Creative Commons Attribution 4.0 International licence, as recorded in the arXiv licence metadata for this exact version and shown on the abstract page https://arxiv.org/abs/2511.00433v1. Full rights notice: licenses/arxiv-2511.00433-CC-BY-4.0.txt.\par\smallskip

\noindent\textit{Editorial scope.} Counted unit: the Theorem whose source label is \texttt{None} in the article (source0.tex lines 178--180, source label \texttt{None}), delivered together with the article's own complete original proof of it (source0.tex lines 181--224, one \texttt{proof} environment of 44 lines). No mathematics is written, repaired or re-derived: statement and proof are the article's own text line for line. Required context retained uncounted: none. Every definition, display and label of those lines is retained unchanged. Omitted from this package and not redistributed: 1--177, 225--622. Nothing inside a retained excerpt is omitted or altered: every retained body is delivered exactly as the article's lines are written, so no omission receipt of any kind accompanies this package. Citations: the retained excerpts carry no citation command at all, so no bibliography entry of the article is transcribed and no reference is redistributed. Reference adaptation: none was needed and none was made; every reference inside the delivered unit resolves inside it, so no unresolved reference or citation is produced. Printed slips kept literal: no printed word, symbol, hypothesis or number of the counted statement or of its proof was altered. Layout and class: the article's own class is \texttt{article} with packages including graphicx, babel, amsthm, amssymb, fancyhdr, pict2e, url, csquotes, amsmath, geometry, tikz, tikz-cd, subcaption, caption; the wrapper is the lane's pinned \texttt{amsart} editorial wrapper at 10pt on A4, which restates the theorem-like environments the retained text needs and adds the article's own macro definitions that the retained text needs (none), with extra wrapper packages (bm, bbm, dsfont, xspace, cancel, mathtools). The delivered numbering is the unit's own. Assets: no retained span contains a figure environment, a graphics-inclusion command, a PDF-page inclusion, a table or a TikZ picture; no image file is copied into this package, the package declares no asset dependency and no image file is redistributed with it. Licence: version arXiv:2511.00433v1 of this article is distributed under the Creative Commons Attribution 4.0 International licence, as recorded in the arXiv licence metadata for this exact version and shown on the abstract page https://arxiv.org/abs/2511.00433v1. A full rights notice is delivered at licenses/arxiv-2511.00433-CC-BY-4.0.txt. Characters: every retained excerpt is pure ASCII, so the delivered engine, XeLaTeX with its default Latin Modern text font, reports no missing character.
\begin{theorem}
    The poset $\text{AutCl}(D_n)$ is a lattice for all positive integer $n$. 
    \end{theorem}

\begin{proof}
     The result holds trivially for $n=1,2$. For $n\ge 3$, let $n=p_1^{t_1}p_2^{t_2}\dots p_k^{t_k}$, where $p_i$'s are distinct primes and $t_i > 0$, for $1\le i \le k$. From the Remark 1, every subgroup of $D_n$ is either of type (1) or of type (2). Let $[H_1]$ and $[H_2]$ be two elements of AutCl$(D_n)$.  To prove AutCl($D_n$) is a lattice, it is sufficient to show that the meet and join of $[H_1]$ and $[H_2]$ exists.\par Consider the following cases: \vspace{0.2cm} \\
\textbf{Case 1:} If both $H_1$ and $H_2$ are of type (1), then $H_1=\bigl<r^{p_1^{u_1}p_2^{u_2}\dots p_k^{u_k}}\bigr>$ and $H_2=\bigl<r^{p_1^{v_1}p_2^{v_2}\dots p_k^{v_k}}\bigr>$, with $0\le u_i, v_i\le t_i$, $1\le i \le k$. We show that,
\begin{equation*}
    [H_1]\vee' [H_2]=[K_1] \hspace{0.2cm} \text{and} \hspace{0.2cm} [H_1]\wedge' [H_2]=[K_2],
\end{equation*}
where 
\begin{equation*}
 K_1=   \bigl<r^{p_1^{\text{min}\{u_1,v_1\}}p_2^{\text{min}\{u_2,v_2\}}\dots p_k^{\text{min}\{u_k,v_k\}}}\bigr> \hspace{0.2cm} \text{and }\hspace{0.2cm} K_2=\bigl<r^{p_1^{\text{max}\{u_1,v_1\}}p_2^{\text{max}\{u_2,v_2\}}\dots p_k^{\text{max}\{u_k,v_k\}}}\bigr>. 
\end{equation*}
Clearly, $[H_1],[H_2]\lesssim[K_1]$ as $H_1, H_2\le K_1$. So, $[K_1]$ is an upper bound of $\{[H_1],[H_2]\}$. In order to show that $[K_1]$ is the least upper bound of $\{[H_1],[H_2]\}$, assume that $[\bar H]$ be an upper bound of $\{[H_1],[H_2]\}$, then there exist $\phi_1,\phi_2 \in \text{Aut}(D_n)$ with $\phi_1(H_1), \phi_2(H_2)\subseteq \bar H$. Moreover, by Theorem 2.3, $[H_1],[H_2]$ are singletons as $H_1, H_2$ are of type (1). Therefore, $\phi_1(H_1)=H_1, \phi_2(H_2)=H_2$ and consequently, $K_1=H_1 \vee H_2 \subseteq \bar H$, which implies $[K_1]\lesssim [\bar H]$.  \par
Clearly, $[K_2]\lesssim [H_1], [H_2]$ as $K_2\le H_1, H_2$ and so, $[K_2]$ is a lower bound of $\{[H_1],[H_2]\}$. To show that $[K_2]$ is the greatest lower bound of $\{[H_1],[H_2]\}$, we assume that $[\widehat H]$ be a lower bound of  $\{[H_1],[H_2]\}$, then there are maps $\phi_1',\phi_2'\in \text{Aut}(D_n)$ with $\phi_1'(\widehat H)\subseteq H_1$ and $\phi_2' (\widehat H)\subseteq H_2$. By Remark 1, $\widehat H$ is of type (1), thus, by Theorem 2.3, $[\widehat H]$ is singleton and $\phi_1'(\widehat H),\phi_2'(\widehat H)=\widehat H$, which implies $\widehat H \le H_1\wedge H_2 =K_2$ and consequently, $[\widehat H]\lesssim [K_2]$. \vspace{0.2cm}\\
\textbf{Case 2:} If $H_1$ is of type (1) and $H_2$ is of type (2) generated by reflections only, then $H_1=$ $\bigl<r^{p_1^{u_1}p_2^{u_2}\dots p_k^{u_k}}\bigr>$ and $H_2=\bigl<r^js\bigr>$, $0\le u_i\le t_i$, $1\le i\le k$ and $0\le j \le n-1$. Note that $[\bigl<r^js\bigr>]=[\bigl<s\bigr>]$, so without the loss of generality, we can choose $H_2=\bigl<s\bigr>$. We show that,
\begin{equation*}
    [H_1]\vee' [H_2]=[K_1] \hspace{0.2cm} \text{and} \hspace{0.2cm} [H_1]\wedge' [H_2]=[K_2],
\end{equation*}
where  
\begin{equation*}
 K_1=   \bigl<r^{p_1^{u_1}p_2^{u_2}\dots p_k^{u_k}},s\bigr> \hspace{0.2cm} \text{and }\hspace{0.2cm} K_2=\{e\}. 
\end{equation*}
Clearly, $[H_1],[H_2]\lesssim [K_1]$ as $H_1,H_2\le K_1$. In order to show that $[K_1]$ is the least upper bound of $\{[H_1],[H_2]\}$, we assume that $[\bar H]$ be an upper bound of $\{[H_1],[H_2]\}$, then there exist $\phi_1,\phi_2 \in \text{Aut}(D_n)$ with $\phi_1(H_1), \phi_2(H_2)\subseteq \bar H$. By Theorem 2.3, $[H_1]$ is singleton, so, $\phi_1(H_1)=H_1\le \bar H$ and $\phi_2(H_2)=$ $ \bigl<r^js\bigr>\le \bar H$, for some $j$, so, $\phi_1(H_1)\vee \phi_2(H_2)\in [K_1]$ and consequently, $[K_1]\lesssim [\bar H]$. \par
Certainly, $[K_2]\lesssim [H_1], [H_2]$ as $K_2=\{e\}$. So $[K_2]$ is a lower bound of $\{[H_1],[H_2]\}$. Let $[\widehat H]$ be a lower bound of $\{[H_1],[H_2]\}$, then $\widehat H$ consists of rotations only, as $[\widehat H]\lesssim [H_1]$, also, $\widehat H$ consists of reflections only, as $[\widehat H]\lesssim [H_2]$. Therefore, $\widehat H=\{e\}$ and hence $[\widehat H]\lesssim [K_2]$. \vspace{0.2cm}\\
\textbf{Case 3}: If $H_1$ is of type (1) and $H_2$ is of type (2), then $H_1=\bigl<r^{p_1^{u_1}p_2^{u_2}\dots p_k^{u_k}}\bigr>$ and $H_2=$$ \bigl<r^{p_1^{v_1}p_2^{v_2}\dots p_k^{v_k}},r^js\bigr>$, $0\le u_i,v_i\le t_i$, $1\le i\le k$ and $0\le j\le p_1^{v_1}p_2^{v_2}\dots p_k^{v_k}-1$. Note that, $ [\bigl<r^{p_1^{v_1}p_2^{v_2}\dots p_k^{v_k}},r^js\bigr>]=[\bigl<r^{p_1^{v_1}p_2^{v_2}\dots p_k^{v_k}},s\bigr>]$, so, without the loss of generality, we can choose $H_2=\bigl<r^{p_1^{v_1}p_2^{v_2}\dots p_k^{v_k}},s\bigr>$. We show that,
\begin{equation*}
    [H_1]\vee' [H_2]=[K_1] \hspace{0.2cm} \text{and} \hspace{0.2cm} [H_1]\wedge' [H_2]=[K_2],
\end{equation*}
where
\begin{equation*}
 K_1=  \bigl <r^{p_1^{\text{min}\{u_1,v_1\}}p_2^{\text{min}\{u_2,v_2\}}\dots p_k^{\text{min}\{u_k,v_k\}}},s\bigr> \hspace{0.2cm} \text{and }\hspace{0.2cm} K_2=\bigl<r^{p_1^{\text{max}\{u_1,v_1\}}p_2^{\text{max}\{u_2,v_2\}}\dots p_k^{\text{max}\{u_k,v_k\}}}\bigr>.
\end{equation*}
Clearly, $[H_1],[H_2]\lesssim [K_1]$ as $H_1,H_2\le K_1$. Let $[\bar H]$ be an upper bound of $\{[H_1],[H_2]\}$, then there are maps $\phi_1,\phi_2\in \text{Aut}(G)$, with $\phi_1(H_1), \phi_2(H_2)\subseteq \bar H$. As $H_1$ is of type (1), by Theorem 2.3, the class $[H_1]$ is singleton, so, $\phi_1(H_1)=H_1\le \bar H$ and also $\phi_2(\bigl<r^{p_1^{v_1}\dots p_k^{v_k}}\bigr>)=\bigl<r^{p_1^{v_1}\dots p_k^{v_k}}\bigr>\le \bar H$. Clearly, $\phi_2(s)\in \bar H$ is a reflection, so, $\phi_1(H_1)\vee\phi_2(H_2)=H_1\vee \bigl<r^{p_1^{v_1}\dots p_k^{v_k}},\phi_2(s)\bigr>\le \bar H$ and hence, $[\phi_1(H_1)\vee \phi_2(H_2)]= [K_1]$, which implies $[K_1]\lesssim [\bar H]$ and consequently, $[K_1]$ is the least upper bound of $\{[H_1],[H_2]\}$. \par
Certainly, $[K_2]\lesssim [H_1], [H_2]$ as $K_2\le H_1,H_2$. So, $[K_2]$ is a lower bound of $\{[H_1],[H_2]\}$. Let $[\widehat H]$ be a lower bound of $\{[H_1],[H_2]\}$, then $\widehat H $ contains rotations only as $[\widehat H]\lesssim [H_1]$, so, $[\widehat H]$ is singleton and consequently, $\widehat H \le H_1, H_2$, which implies $\widehat H \le H_1\wedge H_2=K_2$. So, $[\widehat H]\lesssim [K_2]$ and hence, $[K_2]$ is the greatest lower bound of $\{[H_1],[H_2]\}$.
\vspace{0.2cm} \\ 
\textbf{Case 4:} If both $H_1$ and $H_2$ are of type (2), then without the loss of generality, assume that $H_1=\bigl<r^{p_1^{u_1}p_2^{u_2},\dots p_k^{u_k}},s\bigr>$ and $H_2=\bigl<r^{p_1^{v_1}p_2^{v_2}\dots p_k^{v_k}},s\bigr>$, $0\le u_i,v_i\le t_i$,$1\le i\le k$. Then, we show that,
\begin{equation*}
    [H_1]\vee' [H_2]=[K_1] \hspace{0.2cm} \text{and} \hspace{0.2cm} [H_1]\wedge' [H_2]=[K_2],
\end{equation*}
where 
\begin{equation*}
 K_1=   \bigl<r^{p_1^{\text{min}\{u_1,v_1\}}p_2^{\text{min}\{u_2,v_2\}}\dots p_k^{\text{min}\{u_k,v_k\}}},s\bigr> \hspace{0.2cm} \text{and }\hspace{0.2cm} K_2=\bigl<r^{p_1^{\text{max}\{u_1,v_1\}}p_2^{\text{max}\{u_2,v_2\}}\dots p_k^{\text{max}\{u_k,v_k\}}},s\bigr>. 
\end{equation*}
Clearly, $[H_1],[H_2]\lesssim [K_1]$ as $H_1,H_2\le K_1$. So, $[K_1]$ is an upper bound of $\{[H_1],[H_2]\}$. In order to show that $[K_1]$ is the least upper bound of $\{[H_1],[H_2]\}$, we assume that $[\bar H]$ be an upper bound of $\{[H_1],[H_2]\}$, then $\bigl<r^{p_1^{u_1}\dots p_k^{u_k}}\bigr>,\bigl<r^{p_1^{v_1}\dots p_k^{v_k}}\bigr>\le \bar H$ and also $\bar H$ contains reflections as automorphisms map reflection $s$ to a reflection, say, $r^js$, for some $j$. Therefore, $[K_1]\lesssim [\bar H]$ and hence, $[K_1]$ is the least upper bound of $\{[H_1],[H_2]\}$.\par
Clearly, $[K_2]\lesssim [H_1],[H_2]$ as $K_2\le H_1,H_2$. So, $[K_2]$ is a lower bound of $\{[H_1],[H_2]\}$. To show that $[K_2]$ is the greatest lower bound of $\{[H_1],[H_2]\}$, we assume that $[\widehat H]$ be any lower bound of  $\{[H_1],[H_2]\}$, then there is an automorphic image of $\widehat H$ in $H_1$ and $H_2$. This means that, there is an automorphic image of $\widehat H$ in $K_2$. Thus, by the definition of AutCl($D_n$), $[\widehat H]\lesssim [K_2]$ and consequently, $[K_2]$ is the greatest lower bound of $\{[H_1],[H_2]\}$. 
\end{proof}

\end{document}
